COMPASS formulates prompt robustness through preserving and altering relations
and uses motion-conditioned attribution to define auxiliary training
objectives. The evaluations support improved paraphrase retrieval after
independent text adaptation and dataset-dependent quality gains before that
stage. Their separation is essential: the source-encoder FID and
adapted-encoder retrieval columns do not describe one jointly evaluated system.

The local theory identifies conditions under which representation consistency
can constrain generator sensitivity and output probabilities. Internal
diagnostics show lower measured sensitivity, while revealing that detached
attribution losses have no direct generator-update path in the audited
configuration. The evidence therefore supports checkpoint-specific robustness
rather than causal attribution control. Establishing the broader selective
objective will require matched training comparisons, multiple generation and
training seeds, and direct motion-level evaluation of semantic edits.